\section{Билет 66. Второе начало термодинамики. Энтропия и цикл Карно. КПД идеальной машины Карно с выводом через энтропию. Цикл Карно в координатах pV и TS.}

\begin{definition}
\textbf{Энтропия} $S$ --- функция состояния термодинамической системы, дифференциал которой для обратимого процесса определяется как:
\begin{equation}
    dS = \frac{\delta Q_{\text{обр}}}{T}. \label{eq:entropy_def_66}
\end{equation}
Поскольку $\oint \frac{\delta Q_{\text{обр}}}{T} = 0$, интеграл $\int_1^2 \frac{\delta Q_{\text{обр}}}{T}$ не зависит от пути перехода, а определяется только начальным и конечным состояниями. Это доказывает, что $S$ является функцией состояния.
\end{definition}

\begin{theorem}[КПД цикла Карно через энтропию]
Цикл Карно состоит из двух изотерм ($T_1$, $T_2$) и двух адиабат. На изотермах энтропия изменяется, на адиабатах ($\delta Q=0$) --- постоянна ($S=\text{const}$).
Для полного цикла изменение энтропии равно нулю:
\begin{equation}
    \Delta S_{\text{цикл}} = \oint dS = \frac{Q_1}{T_1} + \frac{Q_2}{T_2} = 0, \label{eq:carnot_entropy}
\end{equation}
где $Q_1 > 0$ --- теплота, полученная от нагревателя, $Q_2 < 0$ --- теплота, отданная холодильнику.
Отсюда следует:
\begin{equation}
    \frac{|Q_2|}{Q_1} = \frac{T_2}{T_1}.
\end{equation}
Подставляя в определение КПД $\eta = 1 - \frac{|Q_2|}{Q_1}$, получаем:
\begin{equation}
    \eta_{\text{К}} = 1 - \frac{T_2}{T_1}. \label{eq:carnot_eff_entropy_66}
\end{equation}
Вывод показывает, что КПД идеальной машины определяется исключительно разностью температур термостатов и не зависит от рабочего тела или конструкции машины.
\end{theorem}

\begin{property}[Цикл Карно в координатах $pV$ и $TS$]
\begin{itemize}
    \item \textbf{Диаграмма $pV$:} Цикл изображается замкнутой кривой, состоящей из двух изотерм ($pV=\text{const}$) и двух адиабат ($pV^\gamma=\text{const}$). Адиабаты идут круче изотерм ($\gamma>1$). Работа газа $A$ численно равна площади фигуры, ограниченной циклом.
    \item \textbf{Диаграмма $TS$:} Цикл Карно представляет собой \textbf{прямоугольник}. Изотермы --- горизонтальные отрезки ($T=\text{const}$), адиабаты --- вертикальные отрезки ($S=\text{const}$, т.к. $\delta Q=0 \Rightarrow dS=0$).
\end{itemize}
\end{property}

\begin{theorem}[Геометрический смысл в $TS$-диаграмме]
В координатах $TS$ площадь под кривой процесса равна количеству теплоты: $Q = \int T \, dS$.
Для цикла Карно:
\begin{align}
    Q_1 &= T_1 (S_B - S_A) = T_1 \Delta S, \\
    |Q_2| &= T_2 (S_C - S_D) = T_2 \Delta S.
\end{align}
Полезная работа равна площади прямоугольника:
\begin{equation}
    A = Q_1 - |Q_2| = (T_1 - T_2) \Delta S.
\end{equation}
КПД выражается как отношение площадей:
\begin{equation}
    \eta = \frac{A}{Q_1} = \frac{(T_1 - T_2) \Delta S}{T_1 \Delta S} = 1 - \frac{T_2}{T_1}.
\end{equation}
$TS$-диаграмма наглядно демонстрирует, что часть полученной теплоты ($T_2 \Delta S$) неизбежно рассеивается в холодильник, что и ограничивает КПД.
\end{theorem}

\begin{remark}
\textbf{Преимущества $TS$-диаграммы:}
\begin{enumerate}
    \item Площадь под кривой напрямую показывает теплоту, а не работу (как в $pV$).
    \item Площадь замкнутого цикла равна полезной работе за цикл.
    \item Обратимые циклы всегда замкнуты, а адиабатные процессы изображаются вертикальными линиями.
    \item Диаграмма широко используется в теплоэнергетике для анализа циклов паровых и газовых турбин, так как позволяет легко оценивать КПД и потери.
\end{enumerate}
\end{remark}

\begin{figure}[htbp]
    \centering
    \begin{tikzpicture}[
        scale=1.2,
        >=Stealth,
        font=\small,
        thick
    ]

    \node[font=\bfseries] at (3.5, 4.5) {Цикл Карно в TS-координатах};

    % === ОСИ ===
    \draw[->] (0, 0) -- (6, 0) node[right, font=\normalsize] {$S$};
    \draw[->] (0, 0) -- (0, 4) node[above, font=\normalsize] {$T$};
    \node[below left] at (0, 0) {$O$};

    % === ПАРАМЕТРЫ ЦИКЛА ===
    \def\SA{1.5}    % Энтропия левой адиабаты
    \def\SB{4.0}    % Энтропия правой адиабаты
    \def\Ttwo{1.2}  % Температура холодильника (нижняя изотерма)
    \def\Tone{3.2}  % Температура нагревателя (верхняя изотерма)

    % === ЗАШТРИХОВАННАЯ ПЛОЩАДЬ ЦИКЛА (работа A) ===
    \fill[blue!20, opacity=0.6] (\SA, \Ttwo) rectangle (\SB, \Tone);

    % === ПРЯМОУГОЛЬНИК ЦИКЛА ===
    \draw[blue, very thick] (\SA, \Ttwo) rectangle (\SB, \Tone);

    % === ВЕРШИНЫ ЦИКЛА ===
    \fill[red] (\SA, \Tone) circle (3pt);
    \node[above left] at (\SA, \Tone) {$A$};

    \fill[red] (\SB, \Tone) circle (3pt);
    \node[above right] at (\SB, \Tone) {$B$};

    \fill[red] (\SB, \Ttwo) circle (3pt);
    \node[below right] at (\SB, \Ttwo) {$C$};

    \fill[red] (\SA, \Ttwo) circle (3pt);
    \node[below left] at (\SA, \Ttwo) {$D$};

    % === СТРЕЛКИ НАПРАВЛЕНИЯ ОБХОДА (по часовой) ===
    \draw[->, blue, thick] (\SA+0.3, \Tone) -- (\SB-0.3, \Tone);
    \draw[->, blue, thick] (\SB, \Tone-0.3) -- (\SB, \Ttwo+0.3);
    \draw[->, blue, thick] (\SB-0.3, \Ttwo) -- (\SA+0.3, \Ttwo);
    \draw[->, blue, thick] (\SA, \Ttwo+0.3) -- (\SA, \Tone-0.3);

    % === ПОДПИСИ ТЕМПЕРАТУР НА ОСИ T ===
    \draw[dashed, gray] (0, \Tone) -- (\SA, \Tone);
    \draw[dashed, gray] (0, \Ttwo) -- (\SA, \Ttwo);
    \node[left] at (0, \Tone) {$T_1$};
    \node[left] at (0, \Ttwo) {$T_2$};

    % === ПОДПИСИ ЭНТРОПИЙ НА ОСИ S ===
    \draw[dashed, gray] (\SA, 0) -- (\SA, \Ttwo);
    \draw[dashed, gray] (\SB, 0) -- (\SB, \Ttwo);
    \node[below] at (\SA, 0) {$S_A$};
    \node[below] at (\SB, 0) {$S_B$};

    % === РАЗМЕРЫ ΔS И ΔT ===
    \draw[<->, purple, thick] (\SA, -0.7) -- (\SB, -0.7)
        node[midway, below] {$\Delta S = S_B - S_A$};

    % === ПОДПИСИ ПРОЦЕССОВ ===
    \node[blue, font=\footnotesize, align=center] at (2.75, \Tone+0.3)
        {Изотерма $T_1$\\(нагрев, $Q_1 > 0$)};
    \node[blue, font=\footnotesize, align=center] at (2.75, \Ttwo-0.4)
        {Изотерма $T_2$\\(охлаждение, $Q_2 < 0$)};
    \node[blue, font=\footnotesize, rotate=90] at (\SA-0.4, 2.2) {Адиабата};
    \node[blue, font=\footnotesize, rotate=90] at (\SB+0.4, 2.2) {Адиабата};

    % === ФОРМУЛЫ В ЦЕНТРЕ ПРЯМОУГОЛЬНИКА ===
    \node[draw, fill=yellow!10, rounded corners, align=center, font=\small, inner sep=5pt] at (7, 2.2)
        {Площадь = Работа:\\[6pt]
         $A = (T_1 - T_2)\Delta S$\\[6pt]
         $\eta = \dfrac{A}{Q_1} = 1 - \dfrac{T_2}{T_1}$};

    % === БЛОК С ФОРМУЛАМИ ДЛЯ Q1 И Q2 ===
    \node[draw, fill=green!10, rounded corners, align=left, font=\small, inner sep=4pt] at (7, 4.5)
        {Теплота:\\[4pt]
        $Q_1 = T_1 \Delta S$\\[4pt]
        $|Q_2| = T_2 \Delta S$};

    \end{tikzpicture}
    \caption{Цикл Карно в TS-координатах. Прямоугольник $ABCD$: изотермы $T_1$ и $T_2$ (горизонтальные линии), адиабаты (вертикальные линии, $S = \text{const}$). Площадь прямоугольника равна работе $A = (T_1 - T_2)\Delta S$. КПД цикла $\eta = 1 - T_2/T_1$.}
    \label{fig:carnot_ts_diagram}
\end{figure}

